\section*{Capítulo \ref{ch:b1:replication-mitosis} --- Replicación del ADN y mitosis}

\begin{solution}{exo:b1:replication-mitosis:1}
Las polimerasas solo añaden \omterm{def:b1:nucleic-acids:nucleotide}{nucleótidos} a un extremo $3'$, de modo que las
hebras nuevas crecen de $5' \to 3'$; los dos moldes son antiparalelos, así
que en una horquilla una hebra nueva puede crecer hacia la horquilla de
forma continua mientras que la otra debe crecer alejándose de ella, y vuelve
a empezar en fragmentos a medida que se expone más molde.
\end{solution}

\begin{solution}{exo:b1:replication-mitosis:2}
Helicasa (desenrolla), \omterm{def:b1:proteins:peptide}{proteínas} de unión a hebra sencilla (mantienen
separadas las hebras), topoisomerasa (libera la torsión por delante),
primasa (cebadores de \omterm{def:b1:nucleic-acids:chain}{ARN}), \omterm{def:b1:replication-mitosis:fork}{ADN polimerasa} (alarga los cebadores y corrige),
una segunda polimerasa (sustituye los cebadores) y ligasa (sella los
fragmentos).
\end{solution}

\begin{solution}{exo:b1:replication-mitosis:3}
G$_1$: $q$, 2n \omterm{def:b1:genomes:genome}{cromosomas} de una \omterm{def:b1:replication-mitosis:mitosis}{cromátida}. G$_2$: $2q$, 2n \omterm{def:b1:genomes:genome}{cromosomas} de
dos \omterm{def:b1:replication-mitosis:mitosis}{cromátidas}. Tras la \omterm{def:b1:replication-mitosis:mitosis}{mitosis}: $q$ por hija, 2n \omterm{def:b1:genomes:genome}{cromosomas} cada una.
\end{solution}

\begin{solution}{exo:b1:replication-mitosis:4}
Profase (condensación, se forma el \omterm{def:b1:replication-mitosis:mitosis}{huso}), prometafase (se rompe la
envoltura, se unen los cinetocoros), metafase (alineamiento en el ecuador),
anafase (las \omterm{def:b1:replication-mitosis:mitosis}{cromátidas} se separan), telofase (se rehacen las envolturas) y
después citocinesis.
\end{solution}

\begin{solution}{exo:b1:replication-mitosis:5}
La mitad del \omterm{def:b1:genomes:genome}{genoma} es síntesis de hebra retardada: $\num{6.4e9}/150 =
\num{4.3e7}$ fragmentos; y otros tantos cebadores, retiradas y ligaciones.
\end{solution}

\begin{solution}{exo:b1:replication-mitosis:6}
$\num{6.4e9}\times 10^{-5} = \num{64000}$; $\times 10^{-7}$: 640;
$\times 10^{-9}$: unas 6 mutaciones por división.
\end{solution}

\begin{solution}{exo:b1:replication-mitosis:7}
$5000/80 = 62$ divisiones. Un tumor debe dividirse sin límite; sin
telomerasa sus \omterm{def:b1:genomes:karyotype}{telómeros} se agotarían y sus \omterm{def:b1:cell-unit-of-life:cell}{células} se detendrían o
morirían, de modo que casi todos los cánceres reactivan la \omterm{def:b1:enzymes:enzyme}{enzima}.
\end{solution}

\begin{solution}{exo:b1:replication-mitosis:8}
Ciclo $= 1\,\mathrm{h}/0.05 = \qty{20}{h}$; S $= 0.30\times 20 =
\qty{6}{h}$.
\end{solution}

\begin{solution}{exo:b1:replication-mitosis:9}
G$_1$: 8 \omterm{def:b1:genomes:genome}{cromosomas}, 8 \omterm{def:b1:replication-mitosis:mitosis}{cromátidas}, $q$. G$_2$: 8, 16, $2q$. Metafase: 8, 16,
$2q$. Hija: 8, 8, $q$.
\end{solution}

\begin{solution}{exo:b1:replication-mitosis:10}
Empiezan rondas nuevas en el origen cada \qty{20}{min}, antes de que
termine la anterior: la replicación se solapa, y una \omterm{def:b1:cell-unit-of-life:cell}{célula} recién nacida ya
tiene \omterm{def:b1:nucleic-acids:chain}{ADN} parcialmente replicado. Al dividirse, el \omterm{def:b1:genomes:genome}{cromosoma} que termina su
ronda se inició \qty{40}{min} antes, una segunda ronda iniciada
\qty{20}{min} antes tiene las horquillas a mitad de camino y una tercera
acaba de empezar: 4 orígenes y 6 horquillas en la \omterm{def:b1:cell-unit-of-life:cell}{célula} que se divide.
\end{solution}

\begin{solution}{exo:b1:replication-mitosis:11}
No se forma \omterm{def:b1:replication-mitosis:mitosis}{huso}: los \omterm{def:b1:genomes:genome}{cromosomas} se condensan pero no pueden alinearse ni
separarse, y la \omterm{def:b1:cell-unit-of-life:cell}{célula} se detiene en el \omterm{prop:b1:replication-mitosis:checkpoints}{punto de control} del \omterm{def:b1:replication-mitosis:mitosis}{huso} en un
estado semejante a la metafase. Si se escabulle de vuelta a la \omterm{def:b1:replication-mitosis:cycle}{interfase} sin
dividirse, queda con 4n \omterm{def:b1:genomes:genome}{cromosomas} (tetraploide). Las \omterm{def:b1:cell-unit-of-life:cell}{células} detenidas con
\omterm{def:b1:genomes:genome}{cromosomas} condensados y separados son exactamente lo que necesita un
\omterm{def:b1:genomes:karyotype}{cariotipo}: el fármaco se aplica para acumularlas.
\end{solution}

\begin{solution}{exo:b1:replication-mitosis:12}
La replicación garantiza que cada \omterm{def:b1:genomes:genome}{cromosoma} pase a ser dos \omterm{def:b1:replication-mitosis:mitosis}{cromátidas}
idénticas, con un error por cada mil millones; el \omterm{def:b1:replication-mitosis:mitosis}{huso} garantiza que cada
hija reciba una de cada par, uniendo cada \omterm{def:b1:genomes:genome}{cromosoma} desde los dos polos y
reteniendo la anafase hasta que todos están bajo tensión. El \omterm{prop:b1:replication-mitosis:checkpoints}{punto de
control} de G$_2$ asegura que la copia haya terminado antes de que empiece el
recuento; el \omterm{prop:b1:replication-mitosis:checkpoints}{punto de control} del \omterm{def:b1:replication-mitosis:mitosis}{huso} asegura que el recuento esté
preparado antes de la separación. Una \omterm{def:b1:cell-unit-of-life:cell}{célula} que reparte mal termina con un
\omterm{def:b1:genomes:genome}{cromosoma} de más o de menos (aneuploidía): una copia correcta hasta la
letra, entregada en la dirección equivocada, lo que mata a casi todas esas
\omterm{def:b1:cell-unit-of-life:cell}{células} y caracteriza a casi todos los cánceres.
\end{solution}

\begin{solution}{pb:b1:replication-mitosis:1}
\textbf{1.} $50\times 28\,800 = \qty{1.44e6}{bp}$.
\textbf{2.} \qty{2.88e6}{bp}.
\textbf{3.} $\num{6.4e9}/\num{2.88e6} = 2220$ orígenes.
\textbf{4.} \qty{2.9}{Mb}, es decir, $2.9\times 10^6\times 0.34\,\mathrm{nm} =
\qty{0.98}{mm}$ de \omterm{def:b1:nucleic-acids:chain}{ADN} entre orígenes.
\textbf{5.} Separación de $\num{6.4e9}/\num{40000} = \qty{160}{kb}$; cada
horquilla recorre unas \qty{80}{kb}, media hora de trabajo.
\textbf{6.} \qty{125}{Mb} por horquilla a \qty{50}{bp/s}: $\num{2.5e6}$ s,
29 días.
\textbf{7.} La velocidad de la horquilla la limitan la química y la
cromatina (veinte veces más lenta que en las bacterias); los orígenes pueden
multiplicarse sin límite, así que la fase S se acorta añadiendo puntos de
partida, no acelerando las horquillas.
\textbf{8.} La síntesis retardada cubre la mitad del \omterm{def:b1:genomes:genome}{genoma}, $\num{3.2e9}$
\omterm{def:b1:nucleic-acids:nucleotide}{nucleótidos}: $\num{2.1e7}$ fragmentos (las dos horquillas de cada origen
tienen cada una su hebra retardada, y las mitades suman el equivalente de un
\omterm{def:b1:genomes:genome}{genoma}).
\textbf{9.} $\num{2.1e7} + 2\times\num{40000} = \num{2.1e7}$ cebadores.
\textbf{10.} Unas $\num{2.1e7}$ ligaciones (una por fragmento) más una por
cada encuentro de horquillas.
\textbf{11.} $2\times\num{6.4e9} = \num{1.28e10}$ \omterm{def:b1:nucleic-acids:nucleotide}{nucleótidos}.
\textbf{12.} $\num{2.56e10}$ \omterm{def:b1:water-small-molecules:atp}{ATP}; a lo largo de \qty{28800}{s},
$\num{8.9e5}$ por segundo.
\textbf{13.} Alrededor del \qty{0.1}{\%} del recambio de \omterm{def:b1:water-small-molecules:atp}{ATP} de la \omterm{def:b1:cell-unit-of-life:cell}{célula}:
la polimerización es barata; lo caro es fabricar los \omterm{def:b1:nucleic-acids:nucleotide}{nucleótidos} y las
\omterm{def:b1:genomes:nucleosome}{histonas}.
\textbf{14.} \num{64000}, 640 y 6,4 errores.
\textbf{15.} $1500\times 10^{-9} = \num{1.5e-6}$: un \omterm{def:b1:genomes:gene}{gen} de cada
setecientos mil por división.
\textbf{16.} $10^{16}\times 6.4 = \num{6.4e16}$ mutaciones: todos los
cambios individuales posibles del \omterm{def:b1:genomes:genome}{genoma} muchas veces, repartidos entre
$10^{13}$ \omterm{def:b1:cell-unit-of-life:cell}{células}; casi todos caen en \omterm{def:b1:nucleic-acids:chain}{ADN} no codificante o en \omterm{def:b1:cell-unit-of-life:cell}{células} que se
desprenden, y una \omterm{def:b1:cell-unit-of-life:cell}{célula} mutada es una entre miles de millones; solo unas
pocas combinaciones dentro de una misma \omterm{def:b1:cell-unit-of-life:cell}{célula} importan.
\textbf{17.} Cien veces ($10^{-7}$): la sucesión de mutaciones que produce
un cáncer, que a $10^{-9}$ tarda décadas en acumularse, se acumula mucho
antes; los defectos hereditarios de la reparación de desapareamientos causan
cáncer de colon en la edad adulta temprana.
\textbf{18.} $\num{2.3e6}$ pares por horquilla a \qty{1000}{bp/s}:
\qty{2300}{s}, \qty{38}{min}.
\textbf{19.} Unas dos rondas solapadas; una \omterm{def:b1:cell-unit-of-life:cell}{célula} recién nacida lleva dos
orígenes (cada uno ya replicado una vez) y horquillas a mitad de camino: su
\omterm{def:b1:genomes:genome}{cromosoma} está parcialmente replicado al nacer.
\textbf{20.} $\num{4.6e6}/1500 = 3070$ fragmentos por ronda, unos 1,3 por
segundo.
\textbf{21.} $\num{4.6e6}\times 10^{-9} = \num{4.6e-3}$ errores por ronda:
alrededor de una hija de cada 200 lleva una mutación nueva.
\textbf{22.} $10^9\times\num{4.6e-3} = \num{4.6e6}$ mutaciones nuevas por
mililitro; un \omterm{def:b1:genomes:gene}{gen} dado de \qty{1}{kb} recibe $\num{4.6e6}\times
1000/\num{4.6e6} = 1000$ impactos.
\textbf{23.} Con mil mutaciones independientes en cualquier \omterm{def:b1:genomes:gene}{gen} dado por
mililitro y por división, todos los cambios individuales posibles ---
incluidos los que confieren resistencia --- ya están presentes en un cultivo
grande antes de aplicar el antibiótico; el fármaco selecciona, no crea.
\textbf{24.} Ser humano: $\num{6.4e9}$ pares, \num{40000} orígenes,
\qty{50}{bp/s}, \qty{8}{h}. Bacteria: $\num{4.6e6}$ pares, un origen,
\qty{1000}{bp/s}, \qty{40}{min}. Un \omterm{def:b1:genomes:genome}{genoma} mil veces mayor, copiado por
horquillas veinte veces más lentas, en solo doce veces el tiempo: la
diferencia son las decenas de miles de orígenes trabajando en paralelo.
\textbf{25.} Unos \num{2200} orígenes como mínimo, si todos se activaran a
la vez; en realidad unos \num{40000}, uno cada \qty{160}{kb}, que se activan
en sucesión de modo que cada horquilla recorre unas \qty{80}{kb}.
\end{solution}
